\chapter{Bounded Concentration and Finite-Class Learning}
\label{sec:p15-bounded-concentration-learning}

An average can estimate its expectation accurately even when individual
observations vary. A concentration bound states how likely a specified
deviation is under explicit sampling and range assumptions. A uniform
bound over a class permits a method to choose among candidates using
those observations.

The probability foundation here is Hoeffding's original 1963 presentation:
Theorem~1 with all three bounds, Theorem~2 and its two-sample corollary,
and the maximal and martingale extension
\cite[pp.~13--23]{Hoeffding1963BoundedSums}.
The learning and validation results are fully derived applications.
Unlike the consistent-learning setting of \cref{sec:p15-occam-learning},
they allow every candidate to make some errors.

\section{Exponential comparison for one bounded variable}

\begin{hbdefinition}{def:p15-conc-bounded-sum}{Ranges and centered sums}
Let \(X_1,\ldots,X_n\) be mutually independent real random variables,
\(n\geq1\), with deterministic finite bounds
\(a_i\leq X_i\leq b_i\) almost surely and \(a_i\leq b_i\).
Put
\[
 \mu_i=EX_i,\quad S_k=\sum_{i=1}^k(X_i-\mu_i),\quad S_0=0,
 \qquad V_n=\sum_{i=1}^n(b_i-a_i)^2.
\]
Zero-width terms are constant almost surely.
The exponential moment at \(\lambda\in\mathbb R\) is
\(E e^{\lambda X}\); boundedness makes it finite.
All logarithms below are natural.
\index{concentration inequality}\index{exponential moment}
\end{hbdefinition}

\begin{hblemma}{thm:p15-conc-exponential}{proved=annex:p15.conc-exponential}{The chord bound and centered exponential bound}
For \(a\leq X\leq b\) almost surely and \(a<b\), with \(\mu=EX\),
\[
 E e^{\lambda X}
   \leq\frac{b-\mu}{b-a}e^{\lambda a}
          +\frac{\mu-a}{b-a}e^{\lambda b}
                  \qquad(\lambda\in\mathbb R).
\]
For all \(a\leq b\), including \(a=b\),
\[
                     E e^{\lambda(X-\mu)}
                        \leq e^{\lambda^2(b-a)^2/8}.
\]
The constant \(1/8\) cannot be reduced uniformly for all bounded
laws and all real \(\lambda\).
For any real \(W\) with finite exponential moment and
\(\lambda>0\),
\[
                         P(W\geq t)\leq e^{-\lambda t}Ee^{\lambda W}.
\]
\end{hblemma}

\section{The complete common-range bound}

\begin{hbdefinition}{def:p15-conc-binary-rate}{Binary divergence and the mean-dependent coefficient}
For \(0<p<1\) and \(0<q<1\), define
\[
 D(q\Vert p)=q\ln(q/p)+(1-q)\ln((1-q)/(1-p)).
\]
At \(q=0,1\), use the continuous extension, with \(0\ln0=0\).
This is binary relative entropy; the symbol does not denote a
distance with the triangle inequality.
For \(0<p<1\), define
\[
 g(p)=
 \begin{cases}
 \displaystyle\frac{\ln((1-p)/p)}{1-2p},&0<p<1/2,\\[4pt]
 \displaystyle\frac{1}{2p(1-p)},&1/2\leq p<1.
 \end{cases}
\]
The formulas agree continuously at \(p=1/2\), where \(g(p)=2\).
\end{hbdefinition}

\begin{hbtheorem}{thm:p15-conc-common-range}{proved=annex:p15.conc-common-range}{All three bounds of Hoeffding's first theorem}
Suppose \(0\leq X_i\leq1\), and put
\(p=n^{-1}\sum_i EX_i\). For \(0<p<1\) and \(0<t<1-p\),
\[
 P(S_n\geq nt)
 \leq e^{-nD(p+t\Vert p)}
 \leq e^{-nt^2g(p)}
 \leq e^{-2nt^2}.
\]
The coefficient \(g(p)\) is the infimum of
\(D(p+t\Vert p)/t^2\) on \(0<t<1-p\).
For \(p<1/2\) it is attained at \(t=1-2p\);
for \(p\geq1/2\) it is the limit as \(t\downarrow0\).

At \(t=1-p\), the first bound extends to \(p^n\);
independent Bernoulli variables of common mean \(p\) attain it.
For \(t>1-p\), the event is impossible.
If \(p=0\) or \(p=1\), \(S_n=0\) almost surely.
Lower-tail bounds follow by applying the theorem to \(1-X_i\).
For a common interval \([a,b]\), \(a<b\), apply it to
\((X_i-a)/(b-a)\), scaling the deviation and mean accordingly.

The optimized exponential bound is not claimed to be the exact
tail probability at every threshold.
Within the exponential method the first bound cannot be improved: for
independent Bernoulli variables of common mean \(p\), the source's
binomial case (2.5), the minimum over \(h>0\) of the right side of its
(1.7), \(e^{-hnt}Ee^{hS_n}\), equals \(e^{-nD(p+t\Vert p)}\) and is
attained exactly at \(h=\lambda_*\); so the first bound, the source's
(2.1), is the best bound that minimizing (1.7) over \(h\) can give under
these hypotheses \cite[p.~15]{Hoeffding1963BoundedSums}. The source
distinguishes this from the least upper bound of the probability
depending only on \(t\), \(n\), the ranges and \(p\)
\cite[p.~14 and n.~3]{Hoeffding1963BoundedSums}.
Both statements, with \(n=1\) as a case where the least upper bound
\(p/(p+t)\) is strictly smaller, are proved as
\cref{thm:p15-w02b-p15-063-exponential-optimality}.
\end{hbtheorem}

\begin{hbtheorem}{thm:p15-conc-common-comparison}{proved=annex:p15.conc-common-comparison}{The original comparison of the three common-range bounds}
For \(n\geq1,\ 0<p<1,\ 0<t<1-p\), write
\[
 A_1=e^{-nD(p+t\Vert p)},\qquad
 A_2=e^{-nt^2g(p)},\qquad A_3=e^{-2nt^2}.
\]
Then \(A_2=A_3\) exactly when \(p=1/2\).
For \(p>1/2\),
\[
 \exp\!\left\{-\frac{nt^3}{3(1-p)(1-p-t)}\right\}
 <\frac{A_1}{A_2}<
 \exp\!\left\{-\frac{(2p-1)nt^3}{6p^2(1-p)^2}\right\}.
\]
For \(p=1/2\),
\[
 \exp\!\left\{-\frac{4nt^4}{3(1-4t^2)}\right\}
 <\frac{A_1}{A_2}=\frac{A_1}{A_3}
 <e^{-4nt^4/3}.
\]
For \(p<1/2\),
\[
       e^{-n(1-2p-t)^2/p}\leq\frac{A_1}{A_2}\leq1,
 \qquad A_1=A_2\ \Longleftrightarrow\ t=1-2p.
\]
For fixed \(p<1/2\), if \(t\downarrow0\) and
\(nt^2\to a\in[0,\infty)\), the ratio tends to
\[
       \exp\!\left\{-a\left(\frac1{2p(1-p)}-g(p)\right)\right\}.
\]
These are comparisons of valid bounds, not ratios between a bound
and an unknown exact sampling probability.
The source's comparison is its Section~3: the bounds as (3.1)--(3.4),
the estimates for \(p>1/2\) and \(p=1/2\) as (3.6)--(3.7), and, for
\(p<1/2\), the identity (3.8) giving equality at \(t=1-2p\), the lower
estimate (3.9) and the small-\(t\) remark that the limit statement above
makes exact \cite[pp.~18--19]{Hoeffding1963BoundedSums}; the source
writes \(\mu\) for \(p\).
\end{hbtheorem}

\section{Unequal ranges, weights and comparisons}

\begin{hbtheorem}{thm:p15-conc-unequal-ranges}{proved=annex:p15.conc-unequal-ranges}{Hoeffding's second theorem and the two-sample comparison}
If \(V_n>0\), then for every \(t>0\),
\[
 P(S_n\geq t)\leq e^{-2t^2/V_n},\qquad
 P(|S_n|\geq t)\leq 2e^{-2t^2/V_n}.
\]
Each bound may be capped at one. If \(V_n=0\), the
corresponding positive-deviation probabilities are zero.
For fixed real weights \(w_i\), replace \(S_n\) by
\(\sum_iw_i(X_i-EX_i)\) and \(V_n\) by
\(\sum_iw_i^2(b_i-a_i)^2\).

If \(Y_1,\ldots,Y_m,Z_1,\ldots,Z_n\), \(m,n\geq1\), are mutually
independent and all lie in \([a,b]\), \(a<b\), put
\(\bar Y=m^{-1}\sum_{i=1}^mY_i\) and
\(\bar Z=n^{-1}\sum_{j=1}^nZ_j\). Then
\[
 P\{(\bar Y-\bar Z)-(E\bar Y-E\bar Z)\geq t\}
 \leq
 \exp\!\left\{-\frac{2t^2}{(b-a)^2(1/m+1/n)}\right\}.
\]
In contrast, two losses evaluated on the same observation need
not be independent. If the observations are independent and each
loss lies in \([0,1]\), the paired loss differences lie in
\([-1,1]\), giving the valid one-sided bound
\(e^{-nt^2/2}\) for their mean deviation.
\end{hbtheorem}

\begin{hbtheorem}{thm:p15-conc-martingale}{proved=annex:p15.conc-martingale}{Conditional centering and a maximum over a finite horizon}
Let \((\mathcal F_i)\) be a filtration and \(D_i\) adapted,
with \(E(D_i\mid\mathcal F_{i-1})=0\) and deterministic
finite bounds \(a_i\leq D_i\leq b_i\).
Set \(M_k=\sum_{i=1}^kD_i\), \(M_0=0\), and
\(V_n=\sum_{i=1}^n(b_i-a_i)^2\).
For \(V_n>0\) and \(t>0\),
\[
 P\{\max_{0\leq k\leq n}M_k\geq t\}\leq e^{-2t^2/V_n},
 \qquad
 P\{\max_{0\leq k\leq n}|M_k|\geq t\}
                      \leq2e^{-2t^2/V_n}.
\]
If \(V_n=0\), all \(M_k=0\) almost surely.

There is also the source's common-range extension: if
\(0\leq X_i\leq1\) are adapted and
\(E(X_i\mid\mathcal F_{i-1})=\mu_i\) for deterministic
\(\mu_i\), then all three bounds in
\cref{thm:p15-conc-common-range} apply to
\(\max_{0\leq k\leq n}\sum_{i=1}^k(X_i-\mu_i)\geq nt\)
on that theorem's stated interior domain.
\end{hbtheorem}

\begin{hbtheorem}{thm:p15-conc-benchmarks}{proved=annex:p15.conc-benchmarks}{The source's moment and normal benchmarks}
For a centered real \(W\) of variance \(v<\infty\), and \(u>0\),
\[
 P(|W|\geq u)\leq\frac{v}{u^2},\qquad
 P(W\geq u)\leq\frac{v}{v+u^2}.
\]
If \(v=0\), both positive-deviation probabilities are zero.
For a standard normal \(Z\), with density
\(\phi(y)=(2\pi)^{-1/2}e^{-y^2/2}\) and CDF \(\Phi\), \(y>0\),
\[
 \frac{y\phi(y)}{1+y^2}\leq\Phi(-y)\leq\frac{\phi(y)}y,
 \qquad \Phi(-y)\leq e^{-y^2/2}.
\]
In particular \(\Phi(-y)\sim\phi(y)/y\) as \(y\to\infty\).
For independent centered finite rows \(\xi_{n,i}\), write
\(v_n=\sum_iE\xi_{n,i}^2>0\). The source's conditions
\[
 v_n\longrightarrow\infty,\qquad
 \frac{\sum_iE|\xi_{n,i}|^3}{v_n^{3/2}}\longrightarrow0
\]
imply, for every fixed real \(y\),
\[
             P\left\{\sum_i\xi_{n,i}\geq y\sqrt{v_n}\right\}
                                  \longrightarrow\Phi(-y).
\]
This is a limiting tail statement; the preceding concentration
bounds hold at each finite sample size.

Without a common upper range, nonnegative independent variables,
with \(\bar X=n^{-1}\sum_{i=1}^nX_i\) and \(E\bar X=\mu>0\), can attain
\(P\{\bar X-\mu\geq t\}=\mu/(\mu+t)\), \(t>0\).
For \(0\leq X_i\leq1\), the following refinement of the
exponential comparison is also valid, for every \(\lambda>0\)
and \(0<t<1-p\):
\[
 P(S_n\geq nt)\leq
 \frac{(1-p+pe^\lambda)^n-1}
      {e^{\lambda n(p+t)}-1}.
\]
At the optimizing \(\lambda_*\) from the common-range proof,
this is strictly smaller than its first exponential bound.
The source compares its bounds with the normal-model value
\(e^{-nt^2/(2\sigma^2)}\) of its (2.16), which is the minimum over \(h\)
of (1.7) when \(\bar X\) is normal, with \(\sigma^2\) the average
variance: the second bound of \cref{thm:p15-conc-common-range} with
\(p\geq1/2\) equals that value in the binomial case, and the first bound
of \cref{thm:p15-conc-unequal-ranges}, at the threshold \(nt\) for
\(S_n\), equals it when every \(X_i\) takes \(a_i\) and \(b_i\) with
probability \(1/2\) each \cite[p.~17]{Hoeffding1963BoundedSums}. Against
the Chebyshev bounds (1.5)--(1.6), the moment benchmarks above, the
source states that its bounds are better except for small \(t\) or small
\(n\), and that for fixed \(t\) they decrease exponentially in \(n\)
\cite[p.~14]{Hoeffding1963BoundedSums}. These coincidences, the
inequalities they are equality cases of, the exact thresholds in
\(nt^2\) at which the exponential bounds fall below the one-sided and
two-sided Chebyshev values, and the Mills-ratio comparison are stated and
proved as
\cref{thm:p15-w02b-p15-063-chebyshev-comparison,thm:p15-w02b-p15-063-normal-benchmark}.
\end{hbtheorem}

\section{Finite-class learning with bounded loss}

\begin{hbdefinition}{def:p15-conc-learning}{Risk and empirical risk without realizability}
Let \((\mathcal Z,\mathcal B,P)\) be a probability space for one
labelled record \(Z\). Let \(\mathcal H=\{h_1,\ldots,h_r\}\),
\(r\geq1\), be fixed before observing data. For each \(h\),
let its measurable loss \(\ell(h,z)\) lie in \([0,1]\).
An independent sample \(Z_1,\ldots,Z_n\), \(n\geq1\), all with law \(P\), gives
\[
 R(h)=E\ell(h,Z),\qquad
 \widehat R_n(h)=\frac1n\sum_{i=1}^n\ell(h,Z_i).
\]
The losses of different hypotheses on one record may be dependent.
An empirical-risk minimizer chooses a member with the smallest
\(\widehat R_n\), breaking ties by the displayed order.
An \(\eta\)-approximate minimizer, \(\eta\geq0\), has empirical
risk at most \(\min_h\widehat R_n(h)+\eta\).

For classification, \(Z=(X,Y)\), predictions are \(h(X)\),
and \(\ell(h,(X,Y))=\mathbf1_{\{h(X)\ne Y\}}\).
The label need not be determined by \(X\), and no candidate
is assumed to have zero population or sample error.
\index{empirical risk minimization}\index{learning!agnostic finite class}
\end{hbdefinition}

\begin{hbtheorem}{thm:p15-conc-finite-learning}{proved=annex:p15.conc-finite-learning}{Uniform risk control and the selected candidate}
For \(0<\delta<1\), set
\[
                  b_n=\sqrt{\frac{\ln(2r/\delta)}{2n}}.
\]
With probability at least \(1-\delta\), every \(h\in\mathcal H\)
satisfies \(|\widehat R_n(h)-R(h)|\leq b_n\).
On the same event, every \(\eta\)-approximate minimizer
\(\widehat h\) satisfies
\[
 R(\widehat h)\leq\min_{h\in\mathcal H}R(h)+2b_n+\eta.
\]
In particular, exact empirical minimization has excess risk
at most \(\varepsilon>0\) with probability at least \(1-\delta\) if
\[
                 n\geq\left\lceil
                     \frac{2\ln(2r/\delta)}{\varepsilon^2}
                              \right\rceil,\qquad n\geq1.
\]
The clipped intervals
\([\,\max\{0,\widehat R_n(h)-b_n\},
       \min\{1,\widehat R_n(h)+b_n\}\,]\)
cover all risks simultaneously on that event, including the risk
of a candidate selected using the same data.
\end{hbtheorem}

\begin{hbtheorem}{thm:p15-conc-weighted-learning}{proved=annex:p15.conc-weighted-learning}{A countable class with predetermined penalties}
For nonempty countable \(\mathcal H\) and \(0<\delta<1\), let \(w_h>0\),
\(\sum_hw_h\leq1\), be fixed before sampling. Define
\[
                  b_h=\sqrt{\frac{\ln(2/(\delta w_h))}{2n}}.
\]
With probability at least \(1-\delta\), all hypotheses satisfy
\(|\widehat R_n(h)-R(h)|\leq b_h\).
If a measurably selected \(\widehat h\) satisfies
\[
 \widehat R_n(\widehat h)+b_{\widehat h}
       \leq\inf_h\{\widehat R_n(h)+b_h\}+\eta,
\]
then on that same event
\[
              R(\widehat h)\leq
                    \inf_h\{R(h)+2b_h\}+\eta.
\]
For every \(\eta>0\), a measurable such selection exists by
enumerating the class, though finding it need not be computationally
effective. For \(\eta=0\), attainment is an additional requirement.
These weights serve a uniform probability allocation; they are
not posterior model probabilities.
\end{hbtheorem}

\section{Training, validation and selection}

\begin{hbtheorem}{thm:p15-conc-validation}{proved=annex:p15.conc-validation}{What independence of a validation sample permits}
Fix \(0<\delta<1\). Let a random training object \(T\), with any specified probability
law, determine \(K\geq1\) candidate losses
\(\ell_j(T,z)\in[0,1]\), jointly measurable.
Take a validation sample of size \(m\geq1\), independent of \(T\),
with independent records of law \(P\).
Define \(R_j(T)=\int\ell_j(T,z)P(dz)\) and validation averages
\(\widehat R_j(T)\). With probability at least \(1-\delta\),
\[
 \max_{1\leq j\leq K}|\widehat R_j(T)-R_j(T)|
                \leq\sqrt{\frac{\ln(2K/\delta)}{2m}}.
\]
Selecting a candidate by minimum validation loss therefore
costs at most twice this radius relative to the best of these
\(K\) candidates in true risk. For evaluating one fixed
training-derived candidate, \(K=1\) suffices.
Candidates depending further on the validation records are
not automatically covered by this \(K\)-candidate statement.
\end{hbtheorem}

\begin{hbtheorem}{thm:p15-conc-noisy-example}{proved=annex:p15.conc-noisy-example}{A complete classification example with no perfect predictor}
Let \(X,Y\) be binary variables with \(P(X=0)=P(X=1)=1/2\), and suppose
\[
       P(Y=1\mid X=0)=1/4,\qquad P(Y=1\mid X=1)=3/4.
\]
For the four predictors \(0,1,x,1-x\), the risks are
\(1/2,1/2,1/4,3/4\), respectively.
The ordered sample
\(((0,0),(0,1),(1,1),(1,1))\) has probability \(27/4096\);
its empirical risks in that order are \(3/4,1/4,1/4,3/4\).
The displayed tie rule selects the constant-one predictor,
whose excess risk is \(1/4\).

Nevertheless the uniform bound applies. For failure probability
\(1/20\) and desired excess risk \(1/10\), the sufficient sample
size is \(\lceil200\ln160\rceil=1016\).
This is a guaranteed sufficient bound, not a claim of necessity.
\end{hbtheorem}

\begin{hbtheorem}{thm:p15-conc-selection-example}{proved=annex:p15.conc-selection-example}{Searching many candidates for a perfect sample fit}
Take records with \(K\) independent fair binary features and a
constant-zero label. Predictor \(h_j\) outputs feature \(j\).
Every predictor has risk \(1/2\). On \(n\) independent records,
the probability that at least one predictor has empirical risk zero is
\[
                          1-(1-2^{-n})^K.
\]
Thus a selected zero-error candidate need not have small true error.
Applying a single-fixed-candidate probability bound to the selected
candidate discards the effect of this search.
\end{hbtheorem}

\section{Dependence and repeated inspection}

\begin{hbtheorem}{thm:p15-conc-dependence-boundaries}{proved=annex:p15.conc-dependence-boundaries}{Pairwise independence and missing range control}
Mutual independence cannot generally be weakened to pairwise independence
in the independent-sum bound. There exist fifteen pairwise independent
fair binary variables for which
\[
        P\!\left\{\sum_{i=1}^{15}(X_i-1/2)\geq15/2\right\}
                         =1/16>e^{-15/2}.
\]
The full finite construction is supplied in the annex.

There is also no distribution-uniform replacement for the bounded
range using only the sample size and existence of a mean.
For every \(n\geq1\), a nonnegative law of mean one can have
an independent \(n\)-sample average equal to zero with probability
at least \(1/2\). Its range increases with \(n\).
\end{hbtheorem}

\begin{hbtheorem}{thm:p15-conc-time-uniform}{proved=annex:p15.conc-time-uniform}{A simultaneous band for every sample size}
For a sequence satisfying the bounded martingale-difference assumptions
above, with known deterministic \(V_n\), and \(0<\delta<1\),
with probability at least \(1-\delta\) one has for every \(n\geq1\)
\[
                 |M_n|\leq
                 \sqrt{\frac{V_n}{2}\ln\frac{2n(n+1)}{\delta}}.
\]
Use radius zero when \(V_n=0\).
For independent identically distributed \(X_i\in[0,1]\) of
mean \(\mu\), this gives simultaneously
\[
        |\bar X_n-\mu|
             \leq\sqrt{\frac{\ln(2n(n+1)/\delta)}{2n}}
                                  \qquad(n\geq1).
\]
The same event covers any almost surely finite data-dependent stopping
time. The ordinary one-\(n\) bound, by itself, does not establish
this simultaneous statement.
\end{hbtheorem}

\section{Choosing and reporting a bound}

\begin{hbmethod}{def:p15-conc-reference-method}{From a sampling model to a risk guarantee}
Identify the random observations, their dependence structure, the
quantity being averaged and a justified range. Choose the probability
statement before calculating its radius: a fixed mean, a difference,
every candidate in a fixed class, a validation-selected candidate, or
every sample size. Include the relevant class size or fixed weights.

Compute the radius and keep its two interpretations explicit:
it controls a sampling probability under the model, and the radius
describes error in estimating a fixed risk or mean. It is not the
posterior probability of a model and does not bound a different future
population without an additional transport argument.
If the displayed probability bound exceeds one, use one; if a risk
interval extends beyond \([0,1]\), clip it to that range.
\end{hbmethod}

The variance-sensitive bounds and the dependent-sampling constructions
in Hoeffding's Theorem~3 and Sections~5--6 are supplied in
\cref{sec:p15-variance-dependent-concentration}, including the
without-replacement convex comparison. This chapter's independent-input
learning examples use their stated sampling law. Confidence sequences
built from these bounds, and their comparison with reused fixed-sample
intervals, Bonferroni and Holm allocation and Wald's boundaries, are in
\cref{sec:p15-w02b-p15-063}. Statistical confidence procedures are
developed in
\cref{sec:p15-hypothesis-tests-confidence,sec:p15-neyman-confidence-constructions}.
